%!TEX root = ../TM3.tex
%Lukas Moschen

\chapter*{Leseliste}
Als Alternative zum Skriptum wird empfohlen weitere Literatur zu lesen.
\emph{Springer-Link} bietet die Möglichkeit diverse Literatur gratis herunterzuladen, sofern Ihr Rechner über eine IP-Adresse der Universität Innsbruck verfügt.
Gehen Sie dazu wie folgt vor:
\begin{enumerate}[noitemsep]
	\item
		Suche Sie einen der Standorte der \emph{Universitäts- und Landesbibliothek Tirol} auf (Standorte: \href{https://www.uibk.ac.at/ulb/index.html.de}{https://www.uibk.ac.at/ulb/index.html.de}).
	\item
		Melden Sie sich auf einen Rechner mit Ihren MCI-Zugangsdaten an.
	\item
		Rufen Sie die Webseite des \emph{Springer-Link} auf: \href{https://link.springer.com}{https://link.springer.com}
	\item
		Suchen Sie die unten stehende Literatur und laden Sie die diese im pdf-Format herunter.
\end{enumerate}

\paragraph{Kinematik}
\begin{itemize}[noitemsep]
	\item 
		Weiterführende Literatur: \textcite[\ppno~5--35 sowie \ppno~121--139]{Gross.2015}.
	\item
		Weitere Übungsbeispiele: \textcite[\ppno~1--28 und \ppno~73--91]{Gross.2012} sowie \textcite[\ppno~257--262]{Hauger.2017}.
\end{itemize}

\paragraph{Kinetik}
\begin{itemize}[noitemsep]
	\item 
		Weiterführende Literatur: \textcite[\ppno~36--64, \ppno~83--98 sowie \ppno~175--187]{Gross.2015}.
	\item
		Weitere Übungsbeispiele: \textcite[\ppno~29--72 und \ppno~93--138]{Gross.2012} sowie \textcite[\ppno~263--266 und \ppno~270--280]{Hauger.2017}.
\end{itemize}

\paragraph{Schwingungslehre}
\begin{itemize}[noitemsep]
	\item 
		Weiterführende Literatur: \textcite[\ppno~221--230 sowie \ppno~238--262]{Gross.2015}.
	\item
		Weitere Übungsbeispiele: \textcite[\ppno~159--183]{Gross.2012} sowie \textcite[\ppno~263--266 und \ppno~281--284]{Hauger.2017}.
\end{itemize}